\section{Lösungstypen und Lösungsmengen}
\label{sec:linear_systems_solution_sets}

Das Gauß-Verfahren hat im vorigen Abschnitt genau eine Lösung geliefert.
Zunächst untersuchen wir geometrisch, welche anderen Lösungsmengen auftreten
können. Anschließend erkennen wir den jeweiligen Fall algebraisch an der
Zeilenstufenform.

\subsection{Drei geometrische Lösungstypen}
\label{subsec:linear_systems_solution_sets_and_geometry}

Die Lösungsmenge eines LGS kann sehr unterschiedlich aussehen. Wir
untersuchen die möglichen Fälle zunächst an $2\times2$-LGS.

In zwei Dimensionen beschreibt jede lineare Gleichung eine Gerade in der
$x_1$-$x_2$-Ebene. Ein Lösungstupel $(x_1,x_2)$ entspricht einem gemeinsamen
Punkt der Geraden.
\begin{itemize}[leftmargin=1.6em,itemsep=1.2em]
  \item
    \begin{minipage}[t]{\linewidth}
    \textbf{LGS hat eine eindeutige Lösung}\par\smallskip
    \noindent
    \begin{minipage}[c]{0.64\linewidth}
      Wir betrachten das LGS
      \[
        \begin{aligned}
          \textcolor{THMBlue}{x_1+x_2}&\textcolor{THMBlue}{=2},\\
          \textcolor{black!60}{x_1-x_2}&\textcolor{black!60}{=0}.
        \end{aligned}
        \quad\Longleftrightarrow\quad
        \left[
          \begin{array}{@{}rr|r@{}}
            \textcolor{THMBlue}{1}&\textcolor{THMBlue}{1}&\textcolor{THMBlue}{2}\\
            \textcolor{black!60}{1}&\textcolor{black!60}{-1}&\textcolor{black!60}{0}
          \end{array}
        \right]
      \]
      Durch Addition beider Gleichungen folgt $2x_1=2$ und damit $x_1=1$.
      Einsetzen in eine der Gleichungen ergibt $x_2=1$.

      \smallskip
      Geometrische Betrachtung: Die Geraden besitzen genau einen
      Schnittpunkt $S=(1,1)$. Somit ist
      \[
        \mathcal{L}=\Cloze{5cm}{\{(1,1)\}}.
      \]
    \end{minipage}\hfill
    \begin{minipage}[c]{0.32\linewidth}
      \centering
      \CachedTikzFigure{content/resources/chapters/linear_systems/tikz/unique_solution_lines.tex}
    \end{minipage}
    \end{minipage}

  \item
    \begin{minipage}[t]{\linewidth}
    \textbf{LGS hat keine Lösung}\par\smallskip
    \noindent
    \begin{minipage}[c]{0.64\linewidth}
      Wir betrachten das LGS
      \[
        \begin{aligned}
          \textcolor{THMBlue}{x_1+x_2}&\textcolor{THMBlue}{=2},\\
          \textcolor{black!60}{x_1+x_2}&\textcolor{black!60}{=1}.
        \end{aligned}
        \quad\Longleftrightarrow\quad
        \left[
          \begin{array}{@{}rr|r@{}}
            \textcolor{THMBlue}{1}&\textcolor{THMBlue}{1}&\textcolor{THMBlue}{2}\\
            \textcolor{black!60}{1}&\textcolor{black!60}{1}&\textcolor{black!60}{1}
          \end{array}
        \right]
      \]
      Die Gleichungen widersprechen sich gegenseitig und können daher
      nicht gleichzeitig erfüllt sein.

      \smallskip
      Geometrische Betrachtung: Die Geraden sind parallel und besitzen
      keinen Schnittpunkt. Daher gilt
      \[
        \mathcal{L}=\Cloze{5cm}{\emptyset}.
      \]
    \end{minipage}\hfill
    \begin{minipage}[c]{0.32\linewidth}
      \centering
      \CachedTikzFigure{content/resources/chapters/linear_systems/tikz/no_solution_lines.tex}
    \end{minipage}
    \end{minipage}

  \item
    \begin{minipage}[t]{\linewidth}
    \textbf{LGS hat $\infty$-viele Lösungen}\par\smallskip
    \noindent
    \begin{minipage}[c]{0.64\linewidth}
      Wir betrachten das LGS
      \[
        \begin{aligned}
          \textcolor{THMBlue}{x_1+x_2}&\textcolor{THMBlue}{=2},\\
          \textcolor{olive!80!black}{2x_1+2x_2}&\textcolor{olive!80!black}{=4}.
        \end{aligned}
        \quad\Longleftrightarrow\quad
        \left[
          \begin{array}{@{}rr|r@{}}
            \textcolor{THMBlue}{1}&\textcolor{THMBlue}{1}&\textcolor{THMBlue}{2}\\
            \textcolor{olive!80!black}{2}&\textcolor{olive!80!black}{2}&\textcolor{olive!80!black}{4}
          \end{array}
        \right]
      \]
      Die zweite Gleichung ist das Doppelte der ersten und stellt daher
      keine zusätzliche Bedingung dar. Jedes Tupel,
      das die erste Gleichung erfüllt, erfüllt auch die zweite.

      \smallskip
      Geometrische Betrachtung: Die Geraden sind identisch und besitzen
      unendlich viele Schnittpunkte. Das bedeutet nicht, dass jedes Tupel aus
      $\R^2$ eine Lösung ist: Lösungen sind genau die Punkte auf der gemeinsamen
      Geraden. Daher gilt
      \[
        \mathcal{L}=\Cloze{7.5cm}{\{(x_1,x_2)\in\R^2\mid x_1+x_2=2\}}.
      \]
    \end{minipage}\hfill
    \begin{minipage}[c]{0.32\linewidth}
      \centering
      \CachedTikzFigure{content/resources/chapters/linear_systems/tikz/infinitely_many_solutions_lines.tex}
    \end{minipage}
    \end{minipage}
\end{itemize}

\vspace{\baselineskip}

Die geometrische Betrachtung macht anschaulich, dass lediglich diese drei Fälle
auftreten können: Zwei Geraden schneiden sich, verlaufen parallel oder liegen
aufeinander.

Die in $2\mathrm{D}$ gewonnenen Erkenntnisse lassen sich auf beliebige
$m\times n$-LGS verallgemeinern.
\begin{TheoremBox}{Anzahl der Lösungen eines LGS}{linear_systems_three_solution_types}
  Jedes $m\times n$-LGS über $\R$ besitzt genau einen der drei Lösungstypen:
  \SolutionTypeCases
  Insbesondere kann ein LGS weder genau zwei noch allgemein endlich viele
  Lösungen größer als eins besitzen.
\end{TheoremBox}

\subsection{Homogene lineare Gleichungssysteme}
\label{subsec:linear_systems_homogeneous_systems}

Ein wichtiger Spezialfall entsteht, wenn sämtliche Absolutglieder eines LGS
null sind.

\begin{DefinitionBox}{Homogenes lineares Gleichungssystem}{linear_systems_homogeneous_system}
  Ein lineares Gleichungssystem heißt \emph{homogen}
  \EnglishTerm{homogeneous}, wenn alle Absolutglieder gleich $\MathHighlight{\text{null}}$ sind:
  \[
    \begin{aligned}
      a_{11}x_1+a_{12}x_2+\dots+a_{1n}x_n&=\MathHighlight{0},\\
      a_{21}x_1+a_{22}x_2+\dots+a_{2n}x_n&=\MathHighlight{0},\\
      &\mathrel{\phantom{=}}\vdots\\
      a_{m1}x_1+a_{m2}x_2+\dots+a_{mn}x_n&=\MathHighlight{0}.
    \end{aligned}
  \]
  Seine erweiterte Koeffizientenmatrix besitzt die Form $[A\mid 0]$.
\end{DefinitionBox}

\begin{TheoremBox}{Triviale Lösung homogener LGS}{linear_systems_homogeneous_trivial_solution}
  Jedes homogene LGS besitzt die \emph{triviale Lösung}
  \EnglishTerm{trivial solution}
  \[
    (x_1,\dots,x_n)=\Cloze{4.5cm}{(0,\dots,0)}.
  \]
  Denn beim Einsetzen dieser Werte wird jede Gleichung zu $0=0$.
\end{TheoremBox}

\begin{KeyPointBox}{Lösungsmengen homogener LGS}{linear_systems_homogeneous_solution_sets}
  Da die triviale Lösung immer vorhanden ist, kann die Lösungsmenge eines
  homogenen LGS nicht leer sein:
  \SolutionTypeCases*
  Ein homogenes LGS besitzt daher entweder genau eine oder unendlich viele
  Lösungen.
\end{KeyPointBox}

\subsection{Rang und Lösungstyp}
\label{subsec:linear_systems_solution_criteria}

Mit Hilfe der Zeilenstufenform kann man ablesen, zu welchem Lösungstyp
ein LGS gehört. Eine zentrale Rolle nimmt dabei der Begriff des Rangs ein.

\begin{DefinitionBox}{Rang einer Matrix}{linear_systems_matrix_rank}
  Der \emph{Rang} \EnglishTerm{rank} einer Matrix ist die Anzahl der
  Nichtnullzeilen in einer Zeilenstufenform, wobei eine Nichtnullzeile
  mindestens einen Eintrag ungleich $0$ enthält. Man schreibt
  $\operatorname{rang}(A)$.
  \begin{center}
    \CachedTikzFigure{content/resources/chapters/linear_systems/tikz/row_echelon_rank.tex}
  \end{center}
\end{DefinitionBox}

Auf diese Weise können bei einem LGS sowohl der Rang der Koeffizientenmatrix
$A$ als auch der Rang der erweiterten Koeffizientenmatrix $[A\mid b]$
berechnet werden.

Elementare Zeilenumformungen verändern den Rang nicht. Obwohl eine Matrix auf
verschiedenen Wegen in Zeilenstufenform gebracht werden kann, ist diese Anzahl
daher eindeutig.

\begin{ExampleBox}{Rang ablesen}{linear_systems_rank_example}
  \begin{enumerate}[label=\alph*),leftmargin=1.8em,itemsep=1.3em]
    \item
      \[
        [A\mid b]\sim
        \left[
          \begin{array}{@{}rrr|r@{}}
            \phantom{-}1&\phantom{-}2&\phantom{-}0&\phantom{-}3\\
            \phantom{-}0&\phantom{-}0&\phantom{-}1&\phantom{-}5
          \end{array}
        \right]
        \quad\Longrightarrow\quad
        \begin{aligned}
          n&=\Cloze{1cm}{3},\\[0.5em]
          \operatorname{rang}(A)&=\Cloze{1cm}{2},\\[0.5em]
          \operatorname{rang}([A\mid b])&=\Cloze{1cm}{2}
        \end{aligned}
      \]

    \item
      \[
        [A\mid b]\sim
        \left[
          \begin{array}{@{}rr|r@{}}
            \phantom{-}1&-1&\phantom{-}3\\
            \phantom{-}0&\phantom{-}1&\phantom{-}5\\
            \phantom{-}0&\phantom{-}0&\phantom{-}1
          \end{array}
        \right]
        \quad\Longrightarrow\quad
        \begin{aligned}
          n&=\Cloze{1cm}{2},\\[0.5em]
          \operatorname{rang}(A)&=\Cloze{1cm}{2},\\[0.5em]
          \operatorname{rang}([A\mid b])&=\Cloze{1cm}{3}
        \end{aligned}
      \]
  \end{enumerate}
\end{ExampleBox}

\begin{TheoremBox}{Lösbarkeits- und Rangkriterium}{linear_systems_rank_criterion}
  Für ein LGS mit $n$ Unbekannten, Koeffizientenmatrix $A$ und erweiterter
  Koeffizientenmatrix $[A\mid b]$ gilt:
  \[
    \begin{array}{rcll}
      \operatorname{rang}(A)
      &<&\operatorname{rang}([A\mid b])
      &\Longleftrightarrow\quad \text{keine Lösung},\\[0.7em]
      \operatorname{rang}(A)
      &=&\operatorname{rang}([A\mid b])=n
      &\Longleftrightarrow\quad \text{genau eine Lösung},\\[0.7em]
      \operatorname{rang}(A)
      &=&\operatorname{rang}([A\mid b])<n
      &\Longleftrightarrow\quad \infty\text{-viele Lösungen}.
    \end{array}
  \]
\end{TheoremBox}

Der Rangvergleich prüft zunächst die \emph{Verträglichkeit}
\EnglishTerm{consistency} der Gleichungen. Sind die Ränge gleich, entscheidet
der Vergleich von $\operatorname{rang}(A)$ mit der Zahl $n$ der Unbekannten,
ob genau eine oder $\infty$-viele Lösungen vorliegen.

\clearpage
\begin{ExampleBox}{Lösungstyp an der Zeilenstufenform erkennen}{linear_systems_rank_solution_types_example}
  \begin{enumerate}[label=\alph*),leftmargin=1.8em,itemsep=1.3em]
    \item
      \[
        [A\mid b]\sim
        \left[
          \begin{array}{@{}rr|r@{}}
            1&2&3\\
            0&0&1\\
            0&0&0
          \end{array}
        \right]
        \quad\Longrightarrow\quad
        \left.
        \begin{aligned}
          n&=\Cloze{1cm}{2},\\[0.5em]
          \operatorname{rang}(A)&=\Cloze{1cm}{1},\\[0.5em]
          \operatorname{rang}([A\mid b])&=\Cloze{1cm}{2}
        \end{aligned}
        \quad\right\}
        \quad\Longrightarrow\quad
        \Cloze{4cm}{\text{keine Lösung}}
      \]

    \item
      \[
        [A\mid b]\sim
        \left[
          \begin{array}{@{}rr|r@{}}
            1&2&3\\
            0&1&5
          \end{array}
        \right]
        \quad\Longrightarrow\quad
        \left.
        \begin{aligned}
          n&=\Cloze{1cm}{2},\\[0.5em]
          \operatorname{rang}(A)&=\Cloze{1cm}{2},\\[0.5em]
          \operatorname{rang}([A\mid b])&=\Cloze{1cm}{2}
        \end{aligned}
        \quad\right\}
        \quad\Longrightarrow\quad
        \Cloze{4cm}{\text{genau eine Lösung}}
      \]

    \item
      \[
        [A\mid b]\sim
        \left[
          \begin{array}{@{}rrr|r@{}}
            1&2&0&3\\
            0&0&1&5
          \end{array}
        \right]
        \quad\Longrightarrow\quad
        \left.
        \begin{aligned}
          n&=\Cloze{1cm}{3},\\[0.5em]
          \operatorname{rang}(A)&=\Cloze{1cm}{2},\\[0.5em]
          \operatorname{rang}([A\mid b])&=\Cloze{1cm}{2}
        \end{aligned}
        \quad\right\}
        \quad\Longrightarrow\quad
        \Cloze{4cm}{\infty\text{-viele Lösungen}}
      \]
  \end{enumerate}
\end{ExampleBox}


\begin{KeyPointBox}{Widerspruchszeile}{linear_systems_inconsistent_row}
  Ein lineare Gleich, die niemals erfüllt werden kann, ist
  \[
    0x_1+0x_2+\dots+0x_n=c
    \quad\Longleftrightarrow\quad
    \left[\begin{array}{@{}cccc|c@{}}
      0&0&\cdots&0&c
    \end{array}\right],
    \qquad c\neq0,
  \]
  und heißt \emph{Widerspruchszeile} \EnglishTerm{inconsistent row}. 
  Für ein LGS in Zeilenstufenform gilt daher
  \[
    \exists\;\text{Widerspruchszeile}
    \quad\Longleftrightarrow\quad
    \operatorname{rang}(A)<\operatorname{rang}([A\mid b])
    \quad\Longleftrightarrow\quad
    \text{keine Lösung}.
  \]
\end{KeyPointBox}


\begin{ExampleBox}{Lösungstyp eines $3\times3$-LGS bestimmen}{linear_systems_inconsistent_three_by_three_example}
  Bringen Sie die erweiterte Koeffizientenmatrix des folgenden LGS in
  Zeilenstufenform und bestimmen Sie den Lösungstyp:
  \[
    \begin{aligned}
      x_1+x_2+x_3&=3,\\
      2x_1+x_2-x_3&=1,\\
      3x_1+2x_2&=5
    \end{aligned}
    \quad\Longleftrightarrow\quad
    [A\mid b]=
    \left[
      \begin{array}{@{}rrr|r@{}}
        1&1& 1&3\\
        2&1&-1&1\\
        3&2& 0&5
      \end{array}
    \right].
  \]

  \begin{ClozeArea}[32]{35}
    \textbf{Vorwärtselimination:}
    \[
      \begin{array}{r@{\;}r@{\qquad}l}
        &\left[\begin{array}{@{}rrr|r@{}}
          1&1&1&3\\
          2&1&-1&1\\
          3&2&0&5
        \end{array}\right]
        &\begin{array}{@{}l@{}}
          {}\\
          \smash{\TransformOp{Z_2-2Z_1}}\\
          \smash{\TransformOp{Z_3-3Z_1}}
        \end{array}\\[1.0em]
        \Longleftrightarrow
        &\left[\begin{array}{@{}rrr|r@{}}
          1&1&1&3\\
          0&-1&-3&-5\\
          0&-1&-3&-4
        \end{array}\right]
        &\begin{array}{@{}l@{}}
          {}\\
          {}\\
          \smash{\TransformOp{Z_3-Z_2}}
        \end{array}\\[1.0em]
        \Longleftrightarrow
        &\left[\begin{array}{@{}rrr|r@{}}
          1&1&1&3\\
          0&-1&-3&-5\\
          \textcolor{red!70!black}{0}&\textcolor{red!70!black}{0}&
          \textcolor{red!70!black}{0}&\textcolor{red!70!black}{1}
        \end{array}\right]
        &\text{Zeilenstufenform}
      \end{array}
    \]
    Die letzte Zeile ist die
    \textcolor{red!70!black}{\emph{Widerspruchszeile}}:
    \[
      \textcolor{red!70!black}{
        0x_1+0x_2+0x_3=1
        \quad\Longleftrightarrow\quad
        0=1.
      }
    \]
    \textbf{Rangvergleich:}
    \[
      \left.
      \begin{aligned}
        n&=3,\\[0.5em]
        \operatorname{rang}(A)&=2,\\[0.5em]
        \operatorname{rang}([A\mid b])&=3
      \end{aligned}
      \quad\right\}
      \quad\Longrightarrow\quad
      \text{keine Lösung}.
    \]
    Damit ist die Lösungsmenge
    \[
      \mathcal L=\varnothing.
    \]
  \end{ClozeArea}
\end{ExampleBox}

\clearpage
\subsection{Unendlich viele Lösungen parametrisieren}
\label{subsec:linear_systems_parametric_solution_sets}

Besitzt ein lösbares LGS $\infty$-viele Lösungen, wird seine Lösungsmenge
mithilfe von Parametern beschrieben.
Das folgende Beispiel zeigt dies anhand eines bereits bekannten LGS.

\begin{ExampleBox}{Eine Lösungsgerade parametrisieren}{linear_systems_parametric_solution}
  Das folgende $2\times2$-LGS haben wir zuvor bereits geometrisch betrachtet.
  Nun untersuchen wir es zusätzlich mithilfe des Rangs und lösen es mit dem
  Gauß-Verfahren:
  \[
    \begin{aligned}
      \textcolor{THMBlue}{x_1+x_2}&\textcolor{THMBlue}{=2},\\
      \textcolor{olive!80!black}{2x_1+2x_2}&\textcolor{olive!80!black}{=4}.
    \end{aligned}
    \quad\Longleftrightarrow\quad
    \left[
      \begin{array}{@{}rr|r@{}}
        \Cloze{0.7cm}{\textcolor{THMBlue}{1}}&
        \Cloze{0.7cm}{\textcolor{THMBlue}{1}}&
        \Cloze{0.7cm}{\textcolor{THMBlue}{2}}\\[0.35em]
        \Cloze{0.7cm}{\textcolor{olive!80!black}{2}}&
        \Cloze{0.7cm}{\textcolor{olive!80!black}{2}}&
        \Cloze{0.7cm}{\textcolor{olive!80!black}{4}}
      \end{array}
    \right].
  \]
  Die Vorwärtselimination liefert die Zeilenstufenform
  \begin{ClozeArea}[32]{8}
    \[
      \begin{array}{r@{\;}r@{\qquad}l}
        &\left[\begin{array}{@{}rr|r@{}}
          \textcolor{THMBlue}{1}&\textcolor{THMBlue}{1}&\textcolor{THMBlue}{2}\\
          \textcolor{olive!80!black}{2}&\textcolor{olive!80!black}{2}&
          \textcolor{olive!80!black}{4}
        \end{array}\right]
        &\begin{array}{@{}l@{}}
          {}\\
          \smash{\TransformOp{Z_2-2Z_1}}
        \end{array}\\[1.0em]
        \Longleftrightarrow
        &\left[\begin{array}{@{}rr|r@{}}
          \textcolor{THMBlue}{1}&\textcolor{THMBlue}{1}&\textcolor{THMBlue}{2}\\
          \textcolor{olive!80!black}{0}&\textcolor{olive!80!black}{0}&
          \textcolor{olive!80!black}{0}
        \end{array}\right]
        &\text{Zeilenstufenform}
      \end{array}
    \]
  \end{ClozeArea}
  Es gilt
  \[
    \left.
    \begin{aligned}
      n&=\Cloze{1cm}{2},\\[0.5em]
      \operatorname{rang}(A)&=\Cloze{1cm}{1},\\[0.5em]
      \operatorname{rang}([A\mid b])&=\Cloze{1cm}{1}
    \end{aligned}
    \quad\right\}
    \quad\Longrightarrow\quad
    \Cloze{4cm}{\infty\text{-viele Lösungen}}.
  \]

  \noindent
  \begin{minipage}[c]{0.62\linewidth}
    Um die Lösungsmenge geschickt anzugeben, führen wir einen Parameter
    $t\in\R$ ein und setzen $x_2=t$ in Zeile 1 ein.
    \[
      \begin{aligned}
        \Cloze{10cm}{
        x_1+x_2=2
        \quad\overset{x_2\,=\,t}{\Longleftrightarrow}\quad x_1+t=2
        \quad\Longleftrightarrow\quad x_1=2-t.
        }
      \end{aligned}
    \]
    Damit lautet die Lösungsmenge
    \[
      \mathcal{L}
      =\Cloze{7.8cm}{
        \left\{
          \begin{pmatrix}2\\0\end{pmatrix}
          +t\begin{pmatrix}-1\\1\end{pmatrix}
          \;\middle|\;t\in\R
        \right\}
      }.
    \]
  \end{minipage}\hfill
  \begin{minipage}[c]{0.34\linewidth}
    \centering
    \CachedTikzFigure{content/resources/chapters/linear_systems/tikz/infinitely_many_solutions_lines.tex}
  \end{minipage}

  \smallskip
  Der Parameter $t$ wirkt wie ein Schieberegler: Für jeden Wert von $t$
  entsteht eine andere Lösung $(2-t,t)$. Durchläuft $t$ alle reellen Zahlen,
  durchläuft der zugehörige Punkt die gesamte Lösungsgerade $x_1+x_2=2$.
\end{ExampleBox}

\begin{KeyPointBox}{Anzahl der freien Variablen}{linear_systems_number_free_variables}
  Variablen für man einen freien Parameter einführt, nennt man \emph{freie Variablen} 
  \EnglishTerm{free variables}. Für ein lösbares LGS mit $n$ Unbekannten ist die Anzahl
  freier Variablen gegeben durch
  \[
    \#\{\text{freie Variablen}\}
    = \#\{\text{benötigte Parameter}\}
    =n-\operatorname{rang}(A).
  \]
  \par\smallskip
  \hfill
  {\footnotesize\color{black!55}
    $\#\{\ldots\}\mathrel{\widehat{=}}\text{Anzahl der Elemente von }\{\ldots\}$.}
\end{KeyPointBox}


Die Anzahl der freien Variablen bestimmt zugleich die Dimension der
Lösungsmenge:
\[
  \begin{aligned}
    0\ \text{freie Variable}
    &\quad\Longleftrightarrow\quad
    \dim(\mathcal L)=0 \quad \text{(eindeutiger Lösungspunkt)} \\[0.5em]
    1\ \text{freie Variable}
    &\quad\Longleftrightarrow\quad
    \dim(\mathcal L)=1 \quad \text{(Lösungsgerade)} \\[0.5em]
    2\ \text{freie Variablen}
    &\quad\Longleftrightarrow\quad
    \dim(\mathcal L)=2 \quad \text{(Lösungsebene)} \\[0.5em]
    3\ \text{freie Variablen}
    &\quad\Longleftrightarrow\quad
    \dim(\mathcal L)=3 \quad \text{(Lösungsvolumen)} \\[0.5em]
    &\hspace{-2.8em}\vdots &&
  \end{aligned}
\]

